\documentclass[11pt]{article}
\usepackage[T1]{fontenc}
\usepackage{lmodern,amsmath,amssymb,amsthm,mathtools,booktabs,array}
\usepackage[margin=1in]{geometry}
\usepackage{microtype}
\usepackage[hidelinks]{hyperref}
\hypersetup{pdftitle={Power-Radical Descent and Arithmetic Derivatives},pdfauthor={ChatGPT},pdfsubject={Partial abc research with scoped Lean proofs}}
\usepackage{fancyhdr}
\pagestyle{fancy}\fancyhf{}\fancyhead[L]{Power-radical descent and arithmetic derivatives}\fancyhead[R]{ChatGPT}\fancyfoot[C]{\thepage}
\setlength{\headheight}{14pt}
\newtheorem{theorem}{Theorem}[section]
\newtheorem{lemma}[theorem]{Lemma}
\newtheorem{proposition}[theorem]{Proposition}
\newtheorem{corollary}[theorem]{Corollary}
\theoremstyle{definition}\newtheorem{definition}[theorem]{Definition}
\theoremstyle{remark}\newtheorem{remark}[theorem]{Remark}
\DeclareMathOperator{\rad}{rad}\DeclareMathOperator{\ord}{ord}\DeclareMathOperator{\lcm}{lcm}
\newcommand{\Z}{\mathbb Z}\newcommand{\N}{\mathbb N}\newcommand{\R}{\mathbb R}
\newcommand{\vp}{v_p}\newcommand{\Qm}{\mathcal Q_m}
\title{Power-Radical Descent, Unbounded First-Appearance Depth,\\and Pairwise Lifting of Arithmetic Derivatives}
\author{ChatGPT}
\date{September 5, 2026}
\begin{document}
\maketitle
\begin{abstract}
We continue an arithmetic research checkpoint for the abc conjecture. For
coprime positive bases we separate the exact change of radical under a power
map into a repeated-lifting factor dividing the exponent, a first-appearance
factor, and an explicit two-adic correction. This yields same-constant descent
criteria for abc-type inequalities and a multiplicative composition law for
the complete budget. We construct infinitely many bases that are not perfect
powers but whose first-appearance factor at the fixed outer exponent three
exceeds one twenty-eighth of the base. This disproves independent sublinear
first-appearance bounds even after common-power normalization; it does not
disprove abc. Independently, pairwise energy minimization improves a local
integer-lifting error for arithmetic derivatives from the sum to at most the
maximum of the normalized exponents. Twenty-nine standalone Lean theorems
have been compiled and axiom-audited, including the all-depth recurrence and
explicit integer transfer cores. General radical identities and the real
proximity theorem have ordinary proofs here, not full formalizations. The
uniform abc estimate remains unproved in this work.
\end{abstract}
\noindent\textbf{Keywords:} abc conjecture; radicals; multiplicative order; valuation lifting; arithmetic derivatives; Lean.

\section{Problem, starting point, and scope}
For a positive integer $N$, put $\rad(N)=\prod_{p\mid N}p$. A primitive
positive triple satisfies $a+b=c$ and $\gcd(a,b)=1$, hence is pairwise coprime.
The target is
\begin{equation}\label{eq:abc}
\forall\varepsilon>0\ \exists K_\varepsilon>0\ \forall(a,b,c),\qquad
c\le K_\varepsilon\rad(abc)^{1+\varepsilon}.
\end{equation}
The constant must be independent of the triple, its prime support, and every
auxiliary base or exponent. None of the unrestricted estimates equivalent to
\eqref{eq:abc} is used as a premise in our unconditional results.

The starting repository snapshot is commit
\texttt{d3f7b33b1115538920cb5fcd851a97f0c3a26a3d} of
\texttt{wyx66624/ABCConjecture}. Its previous transverse-lifting checkpoint
\cite{transverse} controls integer realization of arithmetic derivatives,
but identifies its proposed uniform small-transverse-derivative condition
as equivalent to abc. Its earlier exponent-lifting checkpoint
\cite{multiflag} separates repeated lifting from first appearance for
$x^n-1$. We extend the latter to relative radical budgets, descent between
actual triples, and base changes. An independent energy argument sharpens
local realization without claiming to bound the original scalar defect.

For $m\ge1$ define the cleared defect
\begin{equation}\label{eq:cleared}
\Qm(a,b,c)=\frac{c^m}{\rad(abc)^{m+1}}.
\end{equation}
Uniform boundedness of \eqref{eq:cleared} for every positive integer $m$ is
equivalent to \eqref{eq:abc}. Indeed, abc with $\varepsilon=1/m$ gives the
bound after taking the $m$th power. Conversely, given $\varepsilon>0$, choose
$m$ with $1/m\le\varepsilon$, take the $m$th root, and use
$\rad(abc)\ge1$. The exponent choice depends only on $\varepsilon$.

\section{Valuation lifting and the exact relative radical}
Throughout this section, $x>y\ge1$ and $\gcd(x,y)=1$.

\begin{lemma}[Elementary valuation lifting]\label{lem:lte}
If an odd prime $p$ divides $U-V$ and $p\nmid UV$, then
\[
v_p(U^t-V^t)=v_p(U-V)+v_p(t)\qquad(t\ge1).
\]
The formula also holds with $U+V$ and $U^t+V^t$ when $t$ is odd.
If $U,V$ are odd and $t$ is even, then
\[
v_2(U^t-V^t)=v_2(U-V)+v_2(U+V)+v_2(t)-1.
\]
For odd $t$, $v_2(U^t-V^t)=v_2(U-V)$, and
$v_2(U^t+V^t)=v_2(U+V)$.
\end{lemma}
\begin{proof}
Write $U=V+p^h w$ with $p\nmid Vw$. Raising to a positive exponent prime to
$p$ and expanding shows a unique lowest-valuation term $tV^{t-1}p^h w$,
of valuation $h$. Raising to $p$, the first term has valuation $h+1$;
the middle binomial terms have valuation at least $1+2h>h+1$, and the last
has valuation $ph>h+1$ for odd $p$, except $p=3,h=1$ where it is still
$3>2$. Iteration and the decomposition $t=p^s u$ prove the odd-prime formula.
Using $-V$ proves the odd sum case. For two, odd geometric or alternating
geometric sums have odd length and are odd modulo two. The first squaring
contributes $v_2(U-V)+v_2(U+V)$. Every later factor
$U^{2^j}+V^{2^j}$, $j\ge1$, is $2$ modulo $8$, and contributes exactly one.
An odd remaining exponent contributes zero additional valuation.
\end{proof}

\subsection{Difference powers}
Let
\[
b_n=x^n-y^n,\quad g_n=\frac{x^n-y^n}{x-y},\quad
R_1=\rad(xy(x-y)),\quad R_n=\rad(xy b_n).
\]
The actual lifted triple is $(y^n,b_n,x^n)$ and has radical $R_n$.
A prime divisor of $b_n$ is called \emph{new relative to this base pair}
when it does not divide $x-y$. Every such prime is odd. For a new prime set
\[
d_p=\ord_{\mathbb F_p^\times}(xy^{-1}),\qquad
\alpha_p=v_p(x^{d_p}-y^{d_p}).
\]
The inverse exists because $p\nmid xy$. The order is greater than one,
divides both $n$ and $p-1$, and is prime to $p$. Define
\begin{align}
T_n(x,y)&=\prod_{\substack{p\mid b_n\\p\nmid x-y}}p^{\alpha_p-1},
&L_n(x,y)&=\prod_{p\mid b_n}p^{v_p(n)},\label{eq:TL}\\
E_n(x,y)&=
\begin{cases}
2^{v_2(x+y)-1},&x,y\text{ odd and }n\text{ even},\\
1,&\text{otherwise},
\end{cases}
&B_n(x,y)&=E_nL_nT_n.\label{eq:B}
\end{align}
Empty products are one. Notice that $L_n\mid n$.

\begin{theorem}[Exact relative radical identity]\label{thm:balance}
For every $n\ge1$,
\begin{equation}\label{eq:balance}
R_n B_n=R_1 g_n.
\end{equation}
In particular, $R_1\mid R_n$ and $R_n/R_1$ is the product of the new primes,
each counted once.
\end{theorem}
\begin{proof}
All primes in $x-y$ persist in $b_n$, and neither integer has a prime factor
in common with $xy$. For an odd old prime, Lemma~\ref{lem:lte} gives
$v_p(g_n)=v_p(n)$. For a new prime it gives
\[
v_p(g_n)=v_p(b_n)=\alpha_p+v_p(n).
\]
Dividing $g_n$ by the product of new primes thus leaves
$p^{\alpha_p-1+v_p(n)}$ at a new prime and $p^{v_p(n)}$ at an odd old prime.
If two occurs, it is old; its exponent in this quotient is $v_2(n)$ for odd
$n$ (both are zero) and $v_2(n)+v_2(x+y)-1$ for even $n$. This is exactly
\eqref{eq:B}. These identities exhaust the prime divisors of $g_n$.
\end{proof}

\begin{remark}[A necessary two-adic correction]
At $(x,y,n)=(7,1,2)$ one has $g_n=8$, $R_1=R_n=42$, $L_n=2$, $T_n=1$,
and $E_n=4$. Omitting $E_n$ gives a false radical identity. This correction
cannot be silently included in the factor claimed to divide $n$.
\end{remark}

\subsection{Odd sum powers}
For odd $n$ put
\[
b_n^+=x^n+y^n,\quad g_n^+=\frac{x^n+y^n}{x+y},\quad
R_1^+=\rad(xy(x+y)),\quad R_n^+=\rad(xyb_n^+).
\]
For $p\mid b_n^+$, $p\nmid x+y$, let
$d_p=\ord(-xy^{-1})$. This is an odd divisor of $n$, greater than one.
Use $\alpha_p=v_p(x^{d_p}+y^{d_p})$ and define $T_n^+,L_n^+$ as in
\eqref{eq:TL} with the new seed and sum. Set $B_n^+=L_n^+T_n^+$.

\begin{proposition}\label{prop:sum}
For odd $n$, $L_n^+\mid n$ and
\begin{equation}\label{eq:sumbalance}
R_n^+B_n^+=R_1^+g_n^+.
\end{equation}
\end{proposition}
\begin{proof}
For a new prime, apply the odd-sum part of Lemma~\ref{lem:lte} at $d_p$ and
$n/d_p$, both odd. Its valuation is $\alpha_p+v_p(n)$ because $d_p\mid p-1$.
For an old odd prime, the valuation increment is $v_p(n)$. Two, if present,
is old and its valuation does not change. Counting the new primes once
proves the identity exactly as before.
\end{proof}

\section{Same-constant descent and composition}
Write $\Qm^-(x,y;n)=\Qm(y^n,x^n-y^n,x^n)$ and
$\Qm^+(x,y;n)=\Qm(x^n,y^n,x^n+y^n)$. In this section $m\ge1$.

\begin{theorem}[Exact defect transfer]\label{thm:transfer}
For difference powers,
\begin{equation}\label{eq:minusdefect}
\frac{\Qm^-(x,y;n)}{\Qm^-(x,y;1)}
=\frac{x^{m(n-1)}B_n^{m+1}}{g_n^{m+1}}
\le\frac{B_n^{m+1}}{x^{n-1}}.
\end{equation}
For odd sum powers,
\begin{equation}\label{eq:plusdefect}
\frac{\Qm^+(x,y;n)}{\Qm^+(x,y;1)}
=\frac{(B_n^+)^{m+1}}{g_n^+}.
\end{equation}
Consequently a bound $\Qm\le K$ at the seed transfers with the same constant
if $B_n^{m+1}\le x^{n-1}$ in the difference case, or
$(B_n^+)^{m+1}\le g_n^+$ in the sum case.
\end{theorem}
\begin{proof}
Substitute $R_n/R_1=g_n/B_n$ in the quotient of the two cleared defects.
In the difference case the height ratio is $x^{n-1}$; the geometric sum
$g_n=\sum_{j=0}^{n-1}x^{n-1-j}y^j$ is at least $x^{n-1}$. In the sum case
the height ratio equals $g_n^+$, so its $m$th power cancels all but one of
the $m+1$ radical-ratio powers. Every denominator used here is positive.
\end{proof}

\begin{corollary}[An obstruction in a minimal counterexample]\label{cor:descent}
Fix $K>0$. If a triple with minimal height among those satisfying
$\Qm>K$ has the difference-power form with $n\ge2$, then
\[
B_n^{m+1}>x^{n-1}.
\]
For odd $n$ this implies
\begin{equation}\label{eq:criticalT}
T_n>\frac{x^{(n-1)/(m+1)}}{n}.
\end{equation}
In the odd sum-power form with $n\ge3$, it must satisfy
$(B_n^+)^{m+1}>g_n^+\ge x^{n-1}/2$.
\end{corollary}
\begin{proof}
The seed is a primitive triple of smaller height, so its cleared defect is
at most $K$. Apply Theorem~\ref{thm:transfer} contrapositively. For odd
$n$, $B_n=L_nT_n\le nT_n$. Finally
$g_n^+=(x^n+y^n)/(x+y)\ge x^{n-1}/2$.
\end{proof}

\begin{remark}[What descent does not give]
A smallest counterexample to a fixed inequality is a legitimate conditional
object. Its nonexistence is not proved here. Most triples do not have two
entries with a common exponent greater than one, so these power maps do not
supply a predecessor for every triple. Furthermore, \eqref{eq:criticalT}
is not contradicted by our bounds. The next section constructs genuinely
large first-appearance factors after common-power normalization.
\end{remark}

\begin{proposition}[Complete-budget composition]\label{prop:cocycle}
For difference powers and positive integers $u,v$,
\[
B_{uv}(x,y)=B_u(x,y)B_v(x^u,y^u).
\]
For odd $u,v$ the corresponding sum identity holds.
\end{proposition}
\begin{proof}
In the difference case, $g_{uv}(x,y)=g_u(x,y)g_v(x^u,y^u)$, and the
intermediate radical is both $R_u(x,y)$ and $R_1(x^u,y^u)$. Multiply the
two exact radical identities and cancel the positive intermediate radical.
For sums, both quotients are integral when $u,v$ are odd, and the same
argument applies. This proves composition without assigning a prime's
contribution twice.
\end{proof}

First-appearance depth alone does not obey this composition law. For an odd
prime with order $d$ at $(x,y)$, the order at $(x^u,y^u)$ is
$d/\gcd(d,u)$. If $p\mid x^{uv}-y^{uv}$ but $p\nmid x^u-y^u$, its new
first valuation is
\[
\alpha_p+v_p(u),
\]
because $p\nmid d$. Dividing the two first-appearance products gives the
more precise identity
\begin{equation}\label{eq:rebase}
T_v(x^u,y^u)=\frac{T_{uv}(x,y)}{T_u(x,y)}\,
\prod_{\substack{p\mid x^{uv}-y^{uv}\\p\nmid x^u-y^u}}p^{v_p(u)}.
\end{equation}
The final product divides $u$. All primes in it are odd. For example,
$\ord_7(2)=3$ and $v_7(2^3-1)=1$, whereas for the base $2^{7^j}$ and outer
exponent three the first valuation at seven is $j+1$.

A canonical way to remove this ambiguity is to write
$x=X^s,y=Y^s$, where $s$ is the greatest common divisor of all the positive
prime exponents in $x$ and $y$ (ignore the empty profile for $y=1$).
Unique factorization proves uniqueness and maximality of $s$. The pair
$(X,Y)$ is not a simultaneous nontrivial power. Formula~\eqref{eq:rebase}
keeps track of what this normalization reallocates. The obstruction below
shows that normalization is insufficient to make $T$ sublinear.

\subsection{A uniform estimate on a restricted large-exponent stratum}
\begin{theorem}\label{thm:stratum}
Fix $A\ge0$ and $\varepsilon>0$. For odd
$n\ge 2(1+\varepsilon)/\varepsilon$, suppose that the relevant
first-appearance product satisfies $T_n\le n^A$ (or $T_n^+\le n^A$).
Then the corresponding primitive power triple has height $C$ and radical
$R$ satisfying
\[
C\le K_{\varepsilon,A}R^{1+\varepsilon},
\]
where, with $b=(A+1)(1+\varepsilon)$, one may take
\[
K_{\varepsilon,A}=2^{2+\varepsilon}
\max\left\{1,\left(\frac{2b}{\mathrm e\varepsilon\log 2}\right)^b\right\}.
\]
This constant is independent of $x,y,n$ within the stated stratum.
\end{theorem}
\begin{proof}
The exact balance and $L_n\mid n$ give $R\ge x^{n-1}/n^{A+1}$ for
differences and $R\ge x^{n-1}/(2n^{A+1})$ for sums. In both cases
$C\le2x^n$. Thus
\[
\frac{C}{R^{1+\varepsilon}}
\le2^{2+\varepsilon}n^b x^{1+\varepsilon-\varepsilon n}
\le2^{2+\varepsilon}n^b2^{-\varepsilon n/2}.
\]
The maximum of $t^b e^{-ct}$ over $t>0$ is $(b/(\mathrm e c))^b$;
use $c=\varepsilon\log2/2$.
\end{proof}
The hypothesis on $T$ is not proved for all bases. Even among powers this
theorem leaves all bounded outer exponents with unbounded bases untreated.
Those families cannot be absorbed into a finite exceptional list.

\section{A normalized first-appearance obstruction of linear size}
Let $\phi(X)=X^2+X+1$. We now give a fully explicit infinite construction,
without factoring its growing values.

\begin{lemma}[A polynomial lifting orbit]\label{lem:orbit}
Set
\[
r_0=226,\qquad r_{k+1}=r_k+32\phi(r_k)\quad(k\ge0).
\]
Then $r_k\equiv30\pmod{49}$, $r_k\equiv2\pmod4$, and
\[
v_7(\phi(r_k))=v_7(r_k^3-1)=k+2.
\]
The sequence is strictly increasing, $\ord_7(r_k)=3$, and no $r_k$ is a
perfect power of exponent at least two.
\end{lemma}
\begin{proof}
The exact polynomial identity is
\begin{equation}\label{eq:orbitpoly}
\phi(X+32\phi(X))=\phi(X)(1024X^2+1088X+1057).
\end{equation}
At $X\equiv30\pmod{49}$, $\phi(X)\equiv0\pmod{49}$, so the first
residue is preserved. The multiplier in \eqref{eq:orbitpoly} is $42$ modulo
$49$, so it has seven-adic valuation one. Since
$\phi(226)=49\cdot1047$ and $1047\equiv4\pmod7$, induction proves the
exact valuation claim for $\phi$. The residue modulo four is preserved
because the increment is a multiple of $32$. Hence $v_2(r_k)=1$, which
excludes $r_k=z^e$ for $e\ge2$. Modulo seven the residues of $r_k,r_k^2,r_k^3$
are $2,4,1$, respectively, proving the order statement.
Finally $r_k^3-1=(r_k-1)\phi(r_k)$ and $r_k-1\equiv1\pmod7$.
The increment is strictly positive.
\end{proof}

The orbit grows unnecessarily quickly. A bounded representative retains
its useful divisibility and the normalization condition.
\begin{theorem}[Small normalized representatives]\label{thm:small}
For every $k\ge0$, define
\[
x_k=r_k\bmod(4\cdot7^{k+2}),
\]
where the remainder lies between zero and its modulus. Then
\begin{align}
2\le x_k&<28\cdot7^{k+1},&x_k&\equiv2\pmod4,\quad x_k\equiv2\pmod7,\label{eq:smallsize}\\
7^{k+2}&\mid\phi(x_k),& 7^{k+2}&\le\phi(x_k).\label{eq:smalldiv}
\end{align}
No $x_k$ is a nontrivial perfect power, and $x_k\to\infty$.
\end{theorem}
\begin{proof}
Reduction modulo $4\cdot7^{k+2}$ preserves the residues modulo four and
seven and the polynomial value modulo $7^{k+2}$. The residue modulo four
implies $x_k\ge2$ and excludes nontrivial perfect powers. The positive
integer $\phi(x_k)$ is a multiple of $7^{k+2}$ and is therefore at least
that large. Since $X^2+X+1$ is bounded on every bounded interval of
nonnegative integers, \eqref{eq:smalldiv} implies $x_k\to\infty$.
The sequence of small representatives need not be strictly increasing;
no monotonicity of these remainders is used.
\end{proof}

\begin{corollary}[No independent sublinear first-depth bound]\label{cor:linear}
There are infinitely many integers $x$ that are not nontrivial perfect
powers and satisfy
\begin{equation}\label{eq:lineart}
T_3(x,1)>x/28.
\end{equation}
Consequently, for each $0\le\theta<1$ and each $C>0$, some such $x$ satisfies
$T_3(x,1)>Cx^\theta$. In particular, no bound depending only on the outer
exponent can bound $T_n$ uniformly over normalized bases.
\end{corollary}
\begin{proof}
Seven is new relative to $(x_k,1)$ and has order three. Its first valuation
is at least $k+2$, so $T_3(x_k,1)\ge7^{k+1}>x_k/28$. Since $x_k\to\infty$,
$x_k^{1-\theta}/28\to\infty$, proving all the assertions.
\end{proof}

For example the first eight small representatives are
\[
30,\ 1010,\ 3754,\ 51774,\ 387914,\ 3211490,\ 19682350,\ 65800758.
\]
A separate implementation obtains the same representatives by lifting the
simple root $2$ of $\phi$ modulo seven one digit at a time, and imposing
$x\equiv2\pmod4$. At a root $r$ modulo $M=7^e$, the next digit is
$t\equiv-(\phi(r)/M)(2r+1)^{-1}\pmod7$, and $r$ is replaced by $r+tM$.
This implementation does not use the recurrence \eqref{eq:orbitpoly}.

\begin{remark}[Why this is not an abc counterexample]
The true abc ratio for $(1,x^3-1,x^3)$ also contains
$\rad(x(x-1))$ and all new radical primes in $\phi(x)$. An unbounded
first-appearance factor, even one linear in $x$, does not bound that entire
radical from above strongly enough to contradict \eqref{eq:abc}.
The refuted claim is independent control of $T$, not a coupled estimate
involving the seed's radical or defect. At fixed $\varepsilon>0$, an abc
disproof would require the entire ratio $c/\rad(abc)^{1+\varepsilon}$ to be
unbounded. No such family is constructed here.
\end{remark}

\section{An independent pairwise improvement for arithmetic derivatives}
This section sharpens local realization instead of imposing a new
unproved global estimate. Let
\[
N=\prod_{i=1}^t p_i^{e_i}>1,\quad
k=\gcd(e_1,\ldots,e_t),\quad f_i=e_i/k,\quad S=\sum_i f_i.
\]
For $\lambda\in\Z^t$, define
\[
\partial_\lambda(N)=\sum_i \frac{N}{p_i}e_i\lambda_i.
\]
A target $d$ is \emph{attainable} when it is one of these integer values.

\begin{lemma}[The attainable-value ideal]\label{lem:ideal}
Put $r=\prod_i p_i$ and $\rho=\prod_{p_i\mid f_i}p_i$. The attainable
values are exactly $G_N\Z$, where
\[
G_N=\frac Nr\,k\rho.
\]
\end{lemma}
\begin{proof}
The image is the ideal generated by the integer coefficients
$(N/p_i)e_i$. Remove the common factor $(N/r)k$. At a prime outside $r$, the
remaining coefficients have minimum valuation zero since $\gcd_i f_i=1$.
At a prime $p_j\mid r$, the $j$th coefficient is $f_j(r/p_j)$. If
$p_j\nmid f_j$, the minimum is zero. Otherwise all coefficients are
divisible by $p_j$, while some $f_i$ with $i\ne j$ is prime to $p_j$ and its
coefficient has valuation exactly one. Thus their gcd is exactly $\rho$.
The single-prime case has $f_1=1$ and fits the same formula. Bezout's identity
also proves existence of an integer vector for every multiple of $G_N$.
\end{proof}

Set
\[
s_i=\begin{cases}1,&p_i\nmid f_i,\\1/p_i,&p_i\mid f_i,\end{cases}
\qquad m_i=f_i s_i\in\Z_{>0},\qquad q=\frac{d}{N\Omega(N)},
\quad\Omega(N)=kS.
\]
The $s_i$ are the natural spacings of $\lambda_i/p_i$ after the forced
regular-prime residues have been imposed. Only their displayed values,
not this interpretation, are needed for the proof.

\begin{theorem}[Pair-energy proximity]\label{thm:energy}
For every attainable integer $d$ there is an integer vector with
$\partial_\lambda(N)=d$ and
\begin{equation}\label{eq:kap}
\max_i\left|\frac{\lambda_i}{p_i}-q\right|\le\kappa_N,
\qquad
\kappa_N=\max_i\frac1{2S}\sum_{j\ne i}
\frac{f_j(f_i+f_j)s_i s_j}{\gcd(m_i,m_j)}.
\end{equation}
In particular,
\[
\kappa_N\le\max_i f_i.
\]
For $t=1$ it is zero. If every normalized exponent $f_i$ equals one,
$\kappa_N=(t-1)/t<1$.
\end{theorem}
\begin{proof}
Write $\delta_i=\lambda_i/p_i-q$. The target equation is equivalent to
$\sum_i f_i\delta_i=0$. Among feasible integer vectors, minimize
\[
\mathcal E(\lambda)=\sum_i f_i\delta_i^2.
\]
A feasible vector exists by Lemma~\ref{lem:ideal}; each bounded sublevel
contains only finitely many integer vectors, so the minimum exists.
For $i\ne j$, let $g=\gcd(m_i,m_j)$ and $v=s_i s_j/g$. The two-coordinate
move
\[
(\delta_i,\delta_j)\longmapsto
(\delta_i+f_jv,\delta_j-f_iv)
\]
is a legal integer-weight move: the corresponding increments are
$p_i s_i m_j/g$ and $-p_j s_j m_i/g$, both integers. It preserves the target.
Its energy change, with either sign, is
\[
\pm2f_i f_jv(\delta_i-\delta_j)+f_i f_jv^2(f_i+f_j).
\]
Minimality for both signs gives
$|\delta_i-\delta_j|\le v(f_i+f_j)/2$. Since the weighted mean of the
$\delta_i$ is zero,
\[
|\delta_i|=\left|\frac1S\sum_{j\ne i}f_j(\delta_i-\delta_j)\right|
\le\frac1{2S}\sum_{j\ne i}f_j(f_i+f_j)v,
\]
proving \eqref{eq:kap}. Finally $v\le1$ and
$(f_i+f_j)/2\le\max_i f_i$; sum the weights and divide by $S$.
The two special cases follow by direct substitution.
\end{proof}

\begin{corollary}[A terminating pair-move algorithm]\label{cor:algorithm}
Starting from any feasible vector, repeatedly perform a strictly
energy-decreasing integral multiple of a two-coordinate move whenever
one exists. This process terminates and its output satisfies
\eqref{eq:kap}. No polynomial bound on the number of moves is asserted.
\end{corollary}
\begin{proof}
All visited vectors lie in the finite sublevel set determined by the
initial energy, and strict decrease prevents revisiting a vector.
At termination, the positive and negative unit moves are nondecreasing
for every pair. The proof of Theorem~\ref{thm:energy} only needs these
inequalities, not global optimality.
\end{proof}

The previous local bound was $S=\Omega(N)/k$ \cite{transverse}. The new
bound is at most $\max_i f_i$, and is less than one when all $f_i=1$,
regardless of the number or size of the supporting primes. For actual
finite target choices in the replay we obtain:
\begin{center}
\begin{tabular}{lrrr}\toprule
Profile $N$ & Previous bound $S$ & New envelope $\kappa_N$ & Realized error\\\midrule
$2\cdot3\cdot5\cdot7$ & $4$ & $3/4$ & $419/840$\\
$2\cdot3^{100}$ & $101$ & $50$ & $10150/303$\\
$2^2 3^3$ & $5$ & $1/4$ & $1/5$\\
$2^{64}3^{41}$ & $105$ & $16$ & $736/315$\\\bottomrule
\end{tabular}
\end{center}
The ``realized error'' column is one constructed vector per displayed
profile, not an assertion of the optimal error for every target.

\subsection{Why the global scalar obstruction remains}
The prior global construction first chooses admissible values of
$\partial(a),\partial(b),\partial(c)$ in a one-dimensional aggregate
fibre. When all entries exceed one its period is
\[
\ell=\lcm(h_a,h_b,h_c),\qquad
h_N=\frac{G_N}{\gcd(G_N,N)}\mid k_N.
\]
Replacing its local lifting by Theorem~\ref{thm:energy} changes the overhead
to at most
\[
\frac{\ell}{2\min\{\Omega(a),\Omega(b),\Omega(c)\}}
+\max\{\kappa_a,\kappa_b,\kappa_c\}.
\]
Entries equal to one have zero derivative and are treated without this
three-variable period term. This sharpens realization, but the real
minimum is still
\[
h_0=\frac JR\max\left\{
\frac c{\Omega(a)+\Omega(b)},
\frac b{\Omega(a)+\Omega(c)},
\frac a{\Omega(b)+\Omega(c)}\right\},
\]
where $I\Z$ is the exact Wronskian image, $D=abc/R$, and $J=I/D$.
In particular, the scalar factor $c/R$ has not disappeared.
The uniform small-transverse condition remains at abc strength as established
in the earlier checkpoint; the new proximity estimate is not an independent
proof of that condition.

\section{Proof dependencies, validation, and formal scope}
\subsection{Dependency audit}
The main unconditional chains are
\[
\text{binomial valuation lemma}
\ \Longrightarrow\ \text{exact radical balance}
\ \Longrightarrow\ \text{defect transfer and cocycle};
\]
\[
\begin{gathered}
\text{polynomial orbit identity}
\ \Longrightarrow\ \text{all-depth valuation induction}\\
\Longrightarrow\ \text{small representatives}
\ \Longrightarrow\ T_3>x/28;
\end{gathered}
\]
\[
\text{coefficient gcd and Bezout}
\ \Longrightarrow\ \text{feasible energy minimization}
\ \Longrightarrow\ \text{pairwise proximity envelope}.
\]
None uses abc, Szpiro, Vojta, or a uniform assertion about first appearance.
The large-exponent stratum theorem has the explicit input $T\le n^A$;
it must not be labeled an unrestricted result. The transfer theorem has an
explicit seed bound and budget condition, not an unconditional conclusion
for every triple.

The examples $x=7,y=1,n=2$ and $x_k$ refute, respectively, omission of the
two-adic correction and independent sublinear first-depth control. The
rebasing example explains how ordinary repeated lifting can be misclassified
as first appearance. The derivative proof attacks integer realization,
not the already identified scalar height obstruction.

\subsection{Exact finite checks}
The standard-library program \texttt{computation/replay.py} uses integers
and rational numbers throughout. It rejects optimized Python execution
because assertions are part of the checks. The default run records:
\begin{center}
\begin{tabular}{lr}\toprule
Check & Count\\\midrule
Signed radical identities & 3,553\\
Cleared-defect transfer checks & 14,212\\
Complete-budget composition checks & 405\\
Base-change first-depth checks & 180\\
Exact polynomial-orbit valuations & 13\\
Small representatives, two independent algorithms & 129\\
Local energy lifts on actual factorizations & 7,497\\\bottomrule
\end{tabular}
\end{center}
The power grid covers coprime $1\le y<x\le32$, difference exponents one
through seven, and sum exponents one, three, five, seven. Local profiles
include every $2\le N\le2000$ and 500 generated factorizations, each with
three signed or positive attainable targets. Finite checks validate
implementations; all infinite assertions above have proofs independent
of these counts. No scan is used to prove uniformity.

\subsection{What Lean actually verifies}
\texttt{Lean/PowerDescent.lean} imports only \texttt{Std}. Its 29 named
theorems and 29 axiom queries passed Lean 4.32.2 with this recorded identity:
\begin{center}\small
\begin{tabular}{ll}
Source commit & \texttt{67af725edf41a7b4103eae055e9ee19c6be716c6}\\
Source blob & \texttt{eb665612004cc072255419b6fc5af464f4dc1ffe}\\
Actions run / job & \texttt{33976085784} / \texttt{101332907717}
\end{tabular}\end{center}
Warnings are treated as errors. The reported axiom union consists of
\texttt{propext}, \texttt{Classical.choice}, and \texttt{Quot.sound}.
There is no \texttt{sorryAx}. Two failed development versions were repaired
and are not counted as successful checks.

The formalized infinite statements are genuine universal theorems about
the recurrence, exact divisibility by $7^{k+2}$, residues, exclusion of
all nontrivial perfect powers, and existence of the bounded
representatives for every $k$. In particular the module proves
\begin{gather*}
\forall k\ \exists x,\quad 2\le x<28\cdot7^{k+1},\quad
x\equiv2\pmod4,\quad x\equiv2\pmod7,\\
7^{k+2}\mid x^2+x+1,\quad 7^{k+2}\le x^2+x+1,\quad
\forall z\in\N\ \forall e\ge2,\ x\ne z^e,
\end{gather*}
where in the final clause $z$ is any natural number and $e\ge2$.
The Lean source states the quantifiers without the displayed shorthand.
It also checks abstract integer transfer and composition given an explicit
radical-balance equality, and the polynomial energy-exchange identity.

\emph{Not fully formalized here:} multiplicative-order theory, the general
valuation-to-radical map, all real-power statements, the asymptotic
interpretation of $T_3>x/28$, the real-valued energy-minimum theorem, and the
full abc statement. The root-family's arithmetic certificates are stronger
than finite samples, but they do not make the whole paper kernel-checked.

A checksum-pinned isolated installer and verifier is provided for x86\_64
Ubuntu/WSL. Hosted Ubuntu execution is not access to the user's personal
WSL. Existing project toolchains, dependencies, and verified imports are
not changed. There was no independently running research subagent,
external peer review, or full-repository build in this work.

\section{Remaining mathematical obligations}
Power descent would need a new mechanism jointly controlling first
appearance and the seed's radical or defect. Corollary~\ref{cor:linear}
prevents treating the first factor as uniformly sublinear on normalized
bases independently of that seed information. Descent additionally needs
coverage of triples without two entries sharing a nontrivial exponent;
writing a canonical root decomposition does not create such an exponent.
Neither obligation is proved here.

The derivative route needs an independent height estimate, not merely a
smaller integer-realization error. The FCRT and multi-output route still
needs its uniform endpoint estimate; no new edge exists in a no-face region
solely because the scalar bookkeeping has been sharpened. The IUT
all-place compatibility, geometric uniformity, Pell first-appearance, and
other recorded routes retain their earlier unresolved status. They have
not been refuted as entire approaches, nor newly closed in this supplement.

Thus the strongest completed results of this work are the exact descent
mechanism, the normalized infinite obstruction to independent first-depth
bounds, and the local proximity improvement. The final node
\eqref{eq:abc} is still unresolved. No completed abc proof, disproof, or
Lean closed term for it is asserted.

\begin{thebibliography}{9}
\bibitem{transverse}
ChatGPT, \emph{Quantitative Lifting of Arithmetic Derivatives and the
Transverse Obstruction to abc}, project research checkpoint, September 5,
2026. Repository \texttt{wyx66624/ABCConjecture},
\texttt{research/checkpoints/2026\_09\_05\_transverse\_lifting/}, at the
baseline commit identified in Section 1. Ordinary general proofs and
24 previously checked scoped Lean theorems; not an abc proof.
\bibitem{multiflag}
ChatGPT, \emph{Unitary faces, multi-output transport, and exponent lifting},
project research checkpoint, September 5, 2026. Same repository,
\texttt{research/checkpoints/2026\_09\_05\_multiflag/}.
\bibitem{lean}
Lean project, \emph{Lean 4.32.2}, compiler commit
\texttt{f3b06c705e6c85f5314019d5d3baab0fec5b580c}.
The isolated verifier pins the Linux archive by SHA-256 and records actual
compiler and axiom output. This reference documents the formal checking
implementation, not an additional number-theoretic premise.
\end{thebibliography}
\end{document}
